This study addresses the operational optimization of water supply systems, with the objective of minimizing costs and improving efficiency in the use of electrical energy. The theoretical foundation is based on temporal analyzes of water consumption, usage patterns, statistical correlations and optimization techniques. The model with binary variables adopted takes into account energy tariffs that change over time and consumption policies, aiming to develop operational strategies to reduce energy costs. Data analysis highlighted patterns that can be used to formulate effective management strategies. The optimization implemented indicated the possibility of significant savings. Future plans include expanding the database, sensitivity analyzes and the development of new management policies in order to promote more efficient management of electrical energy.

% Separe as Keywords por ponto e vírgula.
\keywords{Water Resources Management; Energy Efficiency; Costs Optimization; Mathematical Modeling.}
